\clearpage
\section{PROJETO CONCEITUAL}

O projeto conceitual busca meios de realizar as funções necessárias ao atendimento das especificações-meta. A passagem é de requisito para função, de função para princípio e, depois, para uma concepção verificável. Presença de um componente ou pontuação de uma alternativa não equivalem a desempenho demonstrado \citep{back2008,rozenfeld2006}.

Esta seção incorpora a revisão de 15 de outubro de 2026: 23 vínculos de requisitos foram separados de um catálogo de 17 funções, com atributos e restrições preservados. A primeira Pugh foi adequada ao nível funcional e resultou em empate qualitativo. A equipe priorizou M1, proteção direta da peça após desencaixotamento no distribuidor. As combinações morfológicas e as três concepções foram alinhadas a essa fronteira. Os escores anteriores da MIM permanecem históricos e exigem reavaliação sob M1; dimensionamento, interfaces, custo e seleção da concepção continuam em desenvolvimento.

\subsection{Modelagem Funcional}

A função global é \textbf{resguardar a louça} no trecho distribuidor--transporte--descarga--movimentação na obra, preservando a condição de retorno do dispositivo. A louça é o objeto protegido; o produto em desenvolvimento é o sistema de proteção e movimentação. A fronteira muda entre alternativas, mas a finalidade permanece.

\noindent\textbf{Situação atual, V0.} A louça chega acondicionada na embalagem do fabricante. Os relatos de peças trincadas e de dificuldade na descarga sustentam a investigação. Não foram medidos a aceleração transmitida, a taxa de quebra, a perda de rigidez da embalagem do fluxo estudado ou o instante de origem do dano. As referências de 60--90\,g e o cenário econômico são fonte, estimativa ou cálculo conforme sua origem; não são resultados de ensaio da equipe. A V0 organiza o que medir.

Amortecer impacto, imobilizar, portar, expor condição e reter fragmentos descrevem efeitos requeridos. Acesso, liberação, união e separação material apoiam manutenção e fim de vida. Uniformidade, umidade, ciclos e reciclabilidade qualificam os efeitos e o ciclo de vida. Custo é restrição de viabilidade. Não é necessário criar uma função independente para cada requisito.

\subsubsection{Relação entre requisitos, funções e restrições}
\label{sec:requisitos-funcoes}

A revisão de 15 de outubro de 2026 distingue o que o produto deve fazer, o desempenho exigido e as condições que restringem a solução. Requisitos funcionais e de desempenho podem aparecer na mesma especificação: amortecer impacto descreve um efeito; limitar o pico de aceleração estabelece seu desempenho. Condições ambientais, durabilidade, interfaces e custo também integram o conjunto técnico \citep{nasarequisitos2023,nasadecomposicao2023}.

A análise parte do efeito necessário para cumprir cada requisito. Quando esse efeito já existe, o requisito qualifica a função; quando uma ação adicional é necessária, deriva-se uma função de suporte. Quando se trata de um limite econômico ou geométrico, registra-se a restrição e sua aplicação. A especificação-meta permanece verificável em todos os casos. Expressar uma função por verbo e objeto não basta: deve existir ação, efeito ou capacidade necessária, com fronteira e vínculo identificáveis. Funções também abrangem suporte e comunicação, além de transformações de energia ou material.

A aba Função Global conserva 23 vínculos, V01--V23, associados aos 23 RP, e mantém um catálogo separado com 17 funções. As classes técnica, estrutural, ergonômica, interativa, semântica, sintática e adicional caracterizam as funções; não indicam se o requisito é funcional ou não funcional. Uma função pode responder a vários requisitos, e um requisito pode depender de várias funções. Os IDs numéricos preservados no catálogo não formam uma sequência completa, pois identificadores redundantes ou usados para atributos foram descontinuados.

\begin{table}[H]
\caption{Encaminhamento dos requisitos de qualidade, suporte e restrição}
\label{tab:requisito-funcao}
\centering\scriptsize
\begin{tabularx}{\textwidth}{L{2.4cm}YL{3.3cm}}
\toprule\textbf{Requisito}&\textbf{Vínculo ou desdobramento funcional}&\textbf{Critério que permanece}\\\midrule
RP12: uniformidade&Qualifica F01, Amortecer impacto; não cria outra função de amortecimento&Coeficiente de variação do pico $\leq10\%$\\
RP13: umidade&Qualifica F11, Suportar carga, após condicionamento; bloquear água só será função auxiliar se necessário ao princípio&Perda de capacidade $\leq12{,}5\%$ após 48\,h em câmara úmida\\
RP14: volume&Restringe F03/F08 e a geometria carregada. F14, Recolher partes, atua no retorno vazio&Medir carregado e vazio; recolhimento não comprova o limite carregado\\
RP15: manutenção&F15, Prover acesso; F17, Liberar partes; F22, Unir partes, viabilizam a operação humana&Ciclo de limpeza, inspeção e reposição $\leq5$\,min\\
RP18: vida&Qualifica proteção e suporte ao reuso; não é nova transformação nem reaproveitamento entre gerações&$\geq20$ ciclos até falha funcional\\
RP19: reciclagem&F17/F19 viabilizam liberação e separação material; composição e cadeia regional permanecem critérios&$\geq95\%$ da massa com cadeia estabelecida\\
RP20: modularidade&Restringe arquitetura e interfaces; acesso e uniões precisam permitir troca independente&Demonstrar $\geq3$ módulos substituíveis\\
RP23: custo&Restringe todas as funções e soluções; não gera bloco físico autônomo&Teto provisório de R\$\,15/ciclo e comparação econômica completa\\
\bottomrule\end{tabularx}
\fonte{Aplicação da análise funcional aos RP e às metas de trabalho da equipe. Valores ainda sujeitos à validação.}
\end{table}

\begin{table}[H]
\caption{Catálogo de funções e capacidades de suporte}
\label{tab:catalogo-funcoes}
\centering\scriptsize
\begin{tabularx}{\textwidth}{L{0.8cm}L{3.5cm}Y}
\toprule\textbf{ID}&\textbf{Função}&\textbf{Efeito ou capacidade requerida}\\\midrule
F01&Amortecer impacto&Atenuar choque e vibração antes da energia chegar à cerâmica\\
F02&Imobilizar louça&Conservar o posicionamento da carga na fronteira definida\\
F03&Portar louça&Transferir a carga entre posições com interfaces de movimentação\\
F04&Expor condição&Tornar a região crítica da cerâmica acessível à inspeção\\
F05&Prover pega&Disponibilizar interface para a mão enluvada aplicar força\\
F06&Reter fragmentos&Conter fragmentos também durante abertura e retirada\\
F07&Acomodar variação&Adaptar o posicionamento à faixa dimensional sem trocar partes no ajuste\\
F08&Compatibilizar carga&Compatibilizar geometria e interfaces com palete e transporte\\
F09&Estabilizar conjunto&Manter o conjunto carregado apoiado contra tombamento ou movimento indesejado\\
F11&Suportar carga&Transmitir compressão pelo caminho resistente do conjunto\\
F14&Recolher partes&Mudar a configuração vazia para reduzir o volume de retorno\\
F15&Prover acesso&Viabilizar acesso às regiões de limpeza, inspeção e reposição\\
F16&Comunicar uso&Comunicar ação, orientação, fechamento e advertências\\
F17&Liberar partes&Desfazer uniões para abertura, manutenção e fim de vida\\
F19&Separar materiais&Viabilizar frações separáveis para destinação no fim de vida\\
F21&Exibir identidade&Exibir identificador permanente e único da unidade\\
F22&Unir partes&Estabelecer as uniões que mantêm o conjunto montado e contido\\
\bottomrule\end{tabularx}
\fonte{Catálogo da aba Função Global. As funções de suporte não pressupõem componentes ou módulos exclusivos.}
\end{table}

F10 foi reunida em F16; F12 foi vinculada a F01; RP13 passou a qualificar F11. Os antigos F18 e F20 foram substituídos pelos vínculos de durabilidade e arquitetura às funções necessárias. O rótulo F23 deixou de ser ID funcional; RP23 permanece restrição de custo. Na morfológica, Prover acesso explicita o suporte antes denominado Permitir manutenção. Recolher partes explicita a ação no retorno antes denominada Limitar volume. A comunicação de uso acompanha a identidade, pois um serial ou RFID sozinho não orienta fechamento e manuseio.

FAST e estruturas alternativas usam os IDs do catálogo e mantêm atributos nas condições de verificação. A morfológica acrescenta os IDs associados aos princípios existentes. A MIM continua avaliando os mesmos 13 portadores e os mesmos 156 julgamentos; foram atualizados os vínculos funcionais, sem criar peças ou módulos para uniformidade, vida ou custo. Nenhuma dessas correções representa resultado de ensaio ou fechamento da seleção.


\begin{figure}[H]
\centering
\begin{tikzpicture}[x=1cm,y=1cm,>=Stealth,
 f/.style={draw=cinzamedio,rounded corners=2pt,align=center,text width=3cm,minimum height=0.9cm,font=\footnotesize,execute at begin node=\setstretch{1}},
 d/.style={draw=linhacinza,align=left,text width=5.4cm,minimum height=1.1cm,font=\footnotesize,execute at begin node=\setstretch{1}},
 a/.style={->,draw=cinzamedio}]
\node[f,fill=cinzaclaro] (fg) at (1.6,-4.8) {\textbf{Resguardar louça}};
\node[f] (im) at (5.7,0) {Imobilizar louça};
\node[f] (am) at (5.7,-1.6) {Amortecer impacto};
\node[f] (sc) at (5.7,-3.2) {Suportar carga};
\node[f] (po) at (5.7,-4.8) {Portar louça};
\node[f] (ex) at (5.7,-6.4) {Expor condição};
\node[f] (re) at (5.7,-8) {Reter fragmentos};
\node[f] (ma) at (5.7,-9.6) {Prover acesso};
\node[d] (imd) at (11.6,0) {Acomodar variação e preservar o posicionamento sob solicitações.};
\node[d] (amd) at (11.6,-1.6) {Atenuar choque e vibração. Uniformidade e vida qualificam o efeito.};
\node[d] (scd) at (11.6,-3.2) {Transmitir compressão pelo caminho resistente. RP13 qualifica após umidade.};
\node[d] (pod) at (11.6,-4.8) {Prover pega, estabilizar e compatibilizar carga. Recolher partes no retorno vazio.};
\node[d] (exd) at (11.6,-6.4) {Prover acesso protegido à área crítica. Choque registrado não prova integridade.};
\node[d] (red) at (11.6,-8) {Preservar barreira durante transporte, abertura e retirada.};
\node[d] (mad) at (11.6,-9.6) {Liberar partes para acessar limpeza e reposição; reunir após a operação.};
\foreach \n/\nd in {im/imd,am/amd,sc/scd,po/pod,ex/exd,re/red,ma/mad}{
 \draw[a] (fg.east) -- ++(0.3,0) |- (\n.west);
 \draw[a] (\n.east) -- (\nd.west);
}
\node[font=\footnotesize] at (6.7,1) {\textbf{Por quê?} $\longleftarrow$ \hspace{1cm} $\longrightarrow$ \textbf{Como?}};
\end{tikzpicture}
\caption{Relações FAST da função global com seus meios funcionais}
\label{fig:fast}
\fonte{Elaboração da equipe a partir dos requisitos e da síntese funcional de \citet{back2008}.}
\end{figure}

As relações da Figura~\ref{fig:fast} indicam finalidade e meios, não sequência de tarefas. Comunicar uso inclui indicar a ação; exibir identidade apoia o controle da unidade. As funções de suporte à manutenção e ao fim de vida estão discriminadas no catálogo da Tabela~\ref{tab:catalogo-funcoes}. Registrar evento de impacto é função complementar exploratória, OP-01, que não substitui expor condição nem integra automaticamente os 23 requisitos.

\newcommand{\estruturaems}[3]{%
\begin{tikzpicture}[x=1cm,y=1cm,>=Stealth,
 b/.style={draw=cinzamedio,align=center,text width=2.55cm,minimum height=1.1cm,font=\footnotesize,execute at begin node=\setstretch{1}},
 o/.style={draw=linhacinza,align=center,text width=2.9cm,minimum height=1.1cm,font=\footnotesize,execute at begin node=\setstretch{1}},
 a/.style={->,draw=cinzamedio}]
\node[b] (mi) at (1.3,0) {\textbf{M}\par #1};
\node[b] (mp) at (5.1,0) {Acomodar e\par imobilizar};
\node[b] (mt) at (8.9,0) {Portar e\par reter fragmentos};
\node[o] (mo) at (12.9,0) {\textbf{M}\par #2};
\draw[a] (mi)--(mp);\draw[a] (mp)--(mt);\draw[a] (mt)--(mo);
\node[b] (ei) at (1.3,-2.1) {\textbf{E}\par Choque, vibração\par e compressão};
\node[b] (ep) at (5.1,-2.1) {Amortecer e\par suportar carga};
\node[o] (eo) at (8.9,-2.1) {\textbf{E}\par Atenuação pretendida\par a verificar};
\draw[a] (ei)--(ep);\draw[a] (ep)--(eo);\draw[a,dashed] (ep)--(mp);
\node[b] (eh) at (12.9,-2.1) {\textbf{E}\par Esforço humano\par de manuseio};
\draw[a] (eh.north) -- ++(0,0.2) -| (mt.south);
\node[b] (si) at (1.3,-4.2) {\textbf{S}\par Condição desconhecida,\par identidade e orientação};
\node[b] (sp) at (5.1,-4.2) {Expor condição;\par comunicar uso;\par exibir identidade};
\node[o] (so) at (8.9,-4.2) {\textbf{S}\par Condição acessível\par à inspeção};
\node[o] (su) at (12.9,-4.2) {Usuário verifica\par a condição antes\par do aceite};
\draw[a] (si)--(sp);\draw[a] (sp)--(so);\draw[a] (so)--(su);\draw[a,dashed] (mp)--(sp);
\ifnum#3=1
\node[b] (oi) at (1.3,-6.3) {\textbf{E}\par Evento no ponto\par do sensor};
\node[b] (op) at (5.1,-6.3) {OP-01\par Registrar evento\par (opcional)};
\node[o] (os) at (8.9,-6.3) {\textbf{S}\par Limiar excedido\par ou não excedido};
\node[o] (ou) at (12.9,-6.3) {Apoio à triagem\par de inspeção;\par sem diagnóstico};
\draw[a] (oi)--(op);\draw[a] (op)--(os);\draw[a] (os)--(ou);
\fi
\end{tikzpicture}}

\subsubsection{Modelagem Funcional 1}

A M1 coloca a fronteira na louça descoberta: a embalagem original é retirada no distribuidor. A proteção pode atuar diretamente sobre a peça. A fronteira pode favorecer a inspeção, mas introduz abertura com risco e tempo ainda não avaliados. A variedade de formas exige levantamento do mix. A área crítica só será considerada verificável quando a barreira do novo dispositivo efetivamente permitir inspeção.

\begin{figure}[H]
\centering\estruturaems{Louça fora da\par embalagem original}{Louça protegida e\par fragmentos contidos\par como objetivos}{0}
\caption{Estrutura funcional M1: fronteira na peça descoberta}
\fonte{Elaboração da equipe. E = energia; M = material; S = sinal. As saídas são pretendidas.}
\end{figure}

\subsubsection{Modelagem Funcional 2}

A M2 atua sobre o volume na embalagem original. A energia atravessa dispositivo, caixa e acondicionamento interno antes da cerâmica; essa resposta precisa ser medida. Preservar o volume pode reduzir intervenção no distribuidor, mas não resolve inspeção: um visor externo só é útil quando houver interface de observação compatível na embalagem original. Contenção e acesso à área crítica devem ser atendidos simultaneamente.

\begin{figure}[H]
\centering\estruturaems{Louça no volume\par da embalagem original}{Volume protegido e\par retirada segura\par como objetivos}{0}
\caption{Estrutura funcional M2: fronteira no volume embalado}
\fonte{Elaboração da equipe. A conexão entre material e sinal exige interface de inspeção ainda em estudo.}
\end{figure}

\subsubsection{Modelagem Funcional 3}

A M3 mantém a fronteira material e a cadeia de energia da M2. Acrescenta, opcionalmente, registro de evento em paralelo. OP-01 informa se um limiar foi excedido no seu ponto de montagem e pode apoiar triagem de inspeção. Não demonstra presença ou ausência de trinca. A fabricante do indicador também distingue evento e dano possível \citep{spotseeindicador}.

\begin{figure}[H]
\centering\estruturaems{Mesmo volume\par embalado da M2}{Mesmos objetivos\par materiais da M2}{1}
\caption{Estrutura funcional M3: volume embalado e sinal de evento opcional}
\fonte{Elaboração da equipe. OP-01 não substitui a cadeia de inspeção.}
\end{figure}

O registro de impacto foi adiado no informacional e permanece nesta revisão como opção experimental, sujeita a escopo, benefício, calibração e custo. O limite de aceleração na louça (RP1) não é automaticamente o limiar de sensor na estrutura: posição, duração e resposta dinâmica podem diferir. RP4 continua exigindo acesso à condição da cerâmica.

\subsubsection{Seleção da Modelagem}

A planilha mantém os julgamentos iniciais para rastrear a triagem. Os pesos são o grau de importância do cliente no Pareto, derivado do Mudge; não são a importância técnica dos RP no QFD. A escala é $-2$, $-1$, $0$, $+1$, $+2$. A recontagem das 253 comparações confirmou 577 pontos no Mudge. Aplicando os graus às notas originais, os totais são os da Tabela~\ref{tab:pugh-inicial}.

\begin{table}[H]
\centering\small\caption{Triagem inicial preservada como histórico}\label{tab:pugh-inicial}
\begin{tabularx}{\textwidth}{L{2cm}L{2.5cm}Y}
\toprule\textbf{Alternativa}&\textbf{Pontuação}&\textbf{Interpretação atual}\\\midrule
M1&12&Julgamentos iniciais sujeitos a revisão e verificação\\
M2&0&Referência; não é alternativa aprovada\\
M3&40&Maior pontuação inicial; superioridade ainda não comprovada\\
\bottomrule\end{tabularx}
\fonte{Recontagem da aba Pugh. A observação anterior de 42 pontos foi corrigida.}
\end{table}

As notas iniciais atribuíram à M3 vantagens de proteção, manutenção, estabilidade e custo que não decorrem da inclusão de sinal. A revisão anterior, também preservada, misturava cobertura funcional com desempenho e deixou 13 critérios AV para cada alternativa. Essa revisão não foi completada com notas de desempenho estimadas sem base: a comparação atual foi explicitamente delimitada ao nível funcional.

O Pugh compara alternativas a um datum, no mesmo nível de abstração, e orienta a discussão sem transformar o total em decisão absoluta \citep{mitpugh2015}. Nesta primeira seleção, a equipe avalia diferenças entre estruturas funcionais. A confirmação de aceleração, força, volume, vida e custo pertence à comparação das concepções, conforme as soluções e análises disponíveis. Os 23 requisitos, suas metas e seus protocolos permanecem obrigatórios.

A comparação atual usa M2 como datum e a escala qualitativa $-1$, $0$, $+1$, sem ponderação adicional. $+1$ indica vantagem funcional identificável; $-1$, desvantagem identificável; $0$, a mesma cobertura ou o mesmo aspecto descrito. Ausência de base recebe AV, sem ser convertida em zero. Os pesos derivados do Mudge/Pareto e a escala de cinco níveis permanecem no histórico e na segunda matriz. Não foram inventados pesos para os novos critérios funcionais, nem usados pontos do QFD como se fossem preferências adicionais.

\begin{table}[H]
\centering\footnotesize
\caption{Primeira Pugh no nível das estruturas funcionais}\label{tab:pugh-funcional}
\begin{tabularx}{\textwidth}{L{0.9cm}Y*{3}{>{\centering\arraybackslash}p{1.05cm}}}
\toprule\textbf{ID}&\textbf{Aspecto comparado}&\textbf{M1}&\textbf{M2}&\textbf{M3}\\\midrule
CF01&Cobertura de proteção, retenção e movimentação&0&0&0\\
CF02&Barreiras pré-existentes no caminho de inspeção&+1&0&0\\
CF03&Continuidade do acondicionamento no distribuidor&$-1$&0&0\\
CF04&Cobertura da informação obrigatória ao usuário&0&0&0\\
CF05&Cobertura de retorno, manutenção e fim de vida&0&0&0\\\midrule
&Vantagens / desvantagens / equivalências&1/1/3&0/0/5&0/0/5\\
&\textbf{Saldo qualitativo}&\textbf{0}&\textbf{0}&\textbf{0}\\
\bottomrule
\end{tabularx}
\fonte{Julgamentos da equipe a partir das fronteiras, dos fluxos e do catálogo funcional. A planilha registra justificativa e origem por alternativa.}
\end{table}

CF01, CF04 e CF05 receberam zero porque as três estruturas incluem os mesmos efeitos obrigatórios de proteção, informação e suporte ao reuso. Esses zeros não afirmam igualdade de desempenho físico. OP-01 não melhora automaticamente esses efeitos e continua fora da configuração base.

Em CF02, M1 elimina a embalagem de origem do caminho de observação da cerâmica após o preparo no distribuidor. A barreira do novo dispositivo permanece. A nota $+1$ registra uma facilidade potencial de acesso, sem demonstrar a área crítica exigida por RP4. M2 e M3 mantêm a mesma barreira de origem; registrar choque não a torna transparente. Em CF03, retirar a caixa e transferir a peça acrescenta uma operação em M1. A nota $-1$ representa essa intervenção necessária, sem atribuir tempo ou taxa de avaria ainda não medidos.

\noindent\textbf{Decisão de planejamento.} Em 15 de outubro de 2026, a equipe priorizou M1: retirar a embalagem original no distribuidor e desenvolver proteção direta da peça. A comparação termina em empate, e a preferência explícita orienta estudar o acesso mais direto, assumindo a contrapartida do preparo. M1 não foi declarada vencedora numérica. M2 permanece alternativa de contingência; M3 não é adotada na base.

A escolha encerra a definição da direção funcional para prosseguir, condicionada à viabilidade da proteção direta. O preparo deve incluir retirada controlada, conferência inicial, transferência segura, identidade da unidade e destino da embalagem. Definir pressão e área de contato, inspeção com o dispositivo fechado e contenção na abertura. Se essas condições não puderem ser resolvidas, reabrir a fronteira funcional.

A planilha conserva os 23 critérios do cliente e explicita seu encaminhamento à segunda Pugh, onde serão comparadas C1/C2/C3 sob M1. A direção funcional apoia-se na descrição comparativa; os protocolos permanecem como base para analisar e confirmar o desempenho das soluções físicas. A introdução posterior de OP-01 na M1 exigiria revisão de escopo e benefício, mantendo a fronteira direta da peça; isso não corresponde à M3, definida sobre o volume embalado.


\subsection{Matriz TRIZ}

Os quatro conflitos fortes do QFD foram preservados: proteção--ajuste; inspeção--retenção; movimentação--recondicionamento; ajuste--recondicionamento. A tabela explora princípios por analogia funcional com definições publicadas. Não se afirma consulta comprovada a cada célula da matriz clássica 39$\times$39 \citep{dtoceantriz,childstriz}.

\begin{table}[H]
\caption{Encaminhamento dos conflitos fortes}\centering\scriptsize
\begin{tabularx}{\textwidth}{L{1.9cm}L{4.4cm}YY}
\toprule\textbf{Par}&\textbf{Princípios explorados}&\textbf{Aplicação proposta}&\textbf{Verificação}\\\midrule
RP1--RP7&Segmentação, dinamicidade, qualidade local&Ajuste separado da região amortecedora&Aceleração em toda a faixa\\
RP4--RP6&Propriedades ópticas, membranas, intermediário&Barreira transparente e interface efetiva à louça&Inspeção e contenção na retirada\\
RP3--RP15&Segmentação, extração, qualidade local&Desgaste removível e acesso à limpeza&Força e tempo no mesmo conjunto\\
RP7--RP15&Dinamicidade, qualidade local, extração&Ajuste acessível e poucas regiões de sujeira&Amplitude e recondicionamento\\
\bottomrule\end{tabularx}
\fonte{QFD e definições de \citet{dtoceantriz} e \citet{childstriz}.}
\end{table}

Foram corrigidas associações anteriores: parâmetro 32 é facilidade de fabricação, enquanto 37 trata de detectar e medir; princípio 28 não significa qualquer engate rápido; princípio 39 é ambiente inerte, não vedação para fragmentos. Estabilidade da composição não foi equiparada sem justificativa ao ângulo de tombamento. Mapeamentos genéricos ainda incertos permanecem pendentes.

As aplicações são propostas. O conflito só poderá ser considerado resolvido depois de verificar ambas as metas no mesmo conjunto. Para declarar consulta à matriz clássica, registrar versão, um par principal de parâmetros e o cruzamento consultado.

\subsection{Matriz Morfológica}

A morfológica organiza quatro possibilidades por função ou conjunto de funções compatíveis e três caminhos de combinação. A coluna de IDs liga cada linha ao catálogo. P1--P4 identificam opções da própria linha, não princípios TRIZ. A revisão contempla estabilização, recolhimento no retorno e acesso para manutenção. Uniformidade qualifica F01; umidade qualifica F11; vida, reciclabilidade, geometria, troca independente e custo permanecem critérios. Comunicação de uso e identidade são verificadas em conjunto, sem presumir que serial ou RFID orientem o operador.

O ajuste por troca de insertos foi retirado das combinações propostas: RP7 exige amplitude sem troca durante ajuste. Isso não proíbe reposição de desgaste na manutenção. Após a escolha funcional de M1, C1/C2/C3 passaram a atuar sobre a peça descoberta. Cintas, guias e batentes exigem apoios protegidos, controle de tensão e área de contato, evitando concentração de carga na cerâmica. Visor, painel ou sistema óptico integra a barreira do novo dispositivo e deve alcançar a área crítica sem abrir o conjunto. A retirada da embalagem original no distribuidor não dispensa contenção nem comprova a área inspecionável.

\begin{table}[H]
\caption{Direções derivadas da morfológica}\centering\footnotesize
\begin{tabularx}{\textwidth}{L{0.9cm}L{3.2cm}YY}
\toprule\textbf{ID}&\textbf{Direção}&\textbf{Combinação proposta}&\textbf{Pendência principal}\\\midrule
C1&Simplicidade e manutenção&Berço, EPP, cintas sobre apoios protegidos, guias e carrinho&Contato direto, tensão, visor e preparo\\
C2&Ajuste e estrutura&Estrutura telescópica, interface resiliente, painel transparente e rolamento removível&Contato direto, massa, limpeza e retorno\\
C3&Mobilidade no canteiro&Painéis reguláveis com regiões transparentes, elastômeros, rodízios com freio e apoios&Contato direto, desníveis, estabilidade e volume\\
\bottomrule\end{tabularx}
\fonte{Combinações preliminares sob M1. OP-01 fica fora das três configurações base.}
\end{table}

EPP, PEAD, alumínio e elastômeros são propostas de material, não desempenho demonstrado nem materiais atribuídos à patente Rengo. A herança da anterioridade é de função ou princípio documentado. Material, geometria e ciclos precisam de fonte própria ou ensaio.

\subsection{Revisão da Modelagem Funcional e Matriz Morfológica}

A revisão alinhou fronteira e inspeção, preservou retirada segura e distinguiu objetivo de resultado. Sob M1, as condições decisivas são preparo e transferência seguros no distribuidor; contato direto sem dano; área crítica acessível; retenção no transporte e na abertura; ajuste na mesma configuração; força; envelope externo; vida; e custo completo. A inclusão de sinal deve ser julgada como complemento, sem alterar automaticamente a fronteira material escolhida.

Reter 90\% da massa não autoriza declarar ausência de risco: fragmentos pequenos que escapem também precisam ser avaliados. Escolha de material não comprova 7,5\,kN, e duas camadas não garantem menor aceleração. São verificações a executar.

\subsection{Matriz Indicadora de Módulos}
\label{sec:mim}

A MIM foi elaborada para a C1 antes da prioridade dada à fronteira M1. A matriz reúne 12 direcionadores e 13 portadores, com OP-01 fora da base. Os 156 escores e suas somas são preservados como histórico de julgamentos de projeto, sem selecionar C1. A descrição dos portadores foi adaptada para a peça descoberta, mas isso não revalida os escores. A mudança exige reavaliar os motivos de modularização, as interfaces e os agrupamentos antes de adotá-los para M1. A planilha sinaliza essa pendência em todas as relações do registro.

A MIM indica motivos para modularizar e candidatos a agrupamento; interfaces e acoplamentos também precisam ser examinados \citep{erixon1998,blees2009}. A escala adotada é 0 (sem motivo identificado), 1 (fraco), 3 (médio), 5 (forte), conforme a legenda do material-base da disciplina \citep{basemim}. A literatura consultada apresenta forte como 9; a convenção da base foi mantida explicitamente, sem atribuir-lhe a escala clássica.

\begin{table}[H]
\centering\scriptsize\caption{Direcionadores e critérios de julgamento da MIM}\label{tab:mim-drivers}
\begin{tabularx}{\textwidth}{L{0.8cm}L{4.5cm}Y}
\toprule\textbf{ID}&\textbf{Direcionador}&\textbf{Critério}\\\midrule
D01 & Carry-over & Reaproveitar o projeto em futuras gerações; não é reuso por ciclos. \\
D02 & Evolução tecnológica & Isolar partes sujeitas a mudança de tecnologia. \\
D03 & Alterações planejadas de projeto & Separar mudanças previstas no plano de produto. \\
D04 & Especificação técnica & Acomodar diferenças de requisitos entre variantes. \\
D05 & Estilo & Separar alterações de aparência ou identidade de marca. \\
D06 & Unidade comum & Usar a mesma parte física nas variantes atuais. \\
D07 & Processo e organização & Separar processos especializados ou unidades de trabalho. \\
D08 & Testes em separado & Caracterizar o portador antes do ensaio integrado. \\
D09 & Compra de produto pronto & Adquirir solução funcional pronta; matéria-prima não basta. \\
D10 & Manutenção e mantenabilidade & Acessar, limpar e repor partes durante a vida. \\
D11 & Atualização & Melhorar o produto em serviço, além de reparar desgaste. \\
D12 & Reciclagem & Separar materiais e facilitar a destinação. \\
\bottomrule\end{tabularx}
\fonte{Critérios de modularização, sintetizados de \citet{erixon1998}, \citet{blees2009} e da base; aplicação própria.}
\end{table}

Carry-over trata de reaproveitar projeto entre gerações, não de usar a mesma unidade por vinte ciclos. Compra de produto pronto não equivale a comprar polímero ou painel transparente. D03 ficou zerado porque não há roteiro de alterações futuras; D05, porque estilo não integra o recorte. Zero explicita ausência de motivo de modularização e não dispensa função ou requisito.

\begin{table}[H]
\centering\scriptsize\caption{Portadores propostos para C1 sob M1, com julgamentos a reavaliar}\label{tab:mim-portadores}
\begin{tabularx}{\textwidth}{L{0.8cm}L{3.7cm}YL{2.5cm}}
\toprule\textbf{ID}&\textbf{Função}&\textbf{Portador proposto}&\textbf{RP principais}\\\midrule
S01 & Amortecer impacto & Calços de EPP com assento e retenção a definir & RP1, RP12, RP13, RP18 \\
S02 & Imobilizar louça & Cintas com tensão limitada e apoios protegidos na cerâmica & RP2 \\
S03 & Acomodar variação & Guias deslizantes com curso e trava definidos & RP7 \\
S04 & Expor condição & Visor na barreira do próprio dispositivo, sobre peça descoberta & RP4, RP6 \\
S05 & Prover pega & Cavas de pega no corpo estrutural & RP3, RP5 \\
S06 & Portar louça & Interface estrutural para carrinho de obra & RP3, RP8, RP14 \\
S07 & Reter fragmentos & Bandeja e barreira envolvente contínua & RP6 \\
S08 & Suportar carga & Colunas e caminhos de carga do corpo polimérico & RP11, RP13 \\
S09 & Unir e liberar partes & Fechos sobre-centro com retenção secundária & RP17, RP20, RP22 \\
S10 & Comunicar uso e exibir identidade & Marcação de uso e código gravado na estrutura permanente & RP10, RP16, RP21 \\
S11 & Estabilizar conjunto & Geometria de base alargada sem rodízios livres & RP9 \\
S12 & Recolher partes no retorno & Laterais recolhíveis mantendo a base comum & RP14, RP17 \\
S13* & Registrar evento (OP-01) & Estudo de indicador comercial ou registrador & OP-01 \\
\bottomrule\end{tabularx}
\fonte{Função Global, morfológica e C1. S13* é opção de estudo, sem inclusão na referência.}
\end{table}

Manutenção, vida, reciclagem, troca e custo são critérios transversais (RP15, RP18, RP19, RP20 e RP23), examinados sobre os portadores. O catálogo distingue 17 funções: F15 viabiliza acesso e F19 separação material; esses suportes são distribuídos nas interfaces. S10 reúne uso e identidade (F16 e F21). Uniformidade qualifica F01 e umidade qualifica F11. S12 recolhe partes no vazio, sem comprovar RP14 no estado carregado. Sob M1, S04 deve permitir acesso real através da própria barreira, com a peça descoberta e sem depender de S13. Apoios e cintas exigem avaliação de pressão sobre a cerâmica. Os escores herdados não demonstram que essas novas interfaces estão resolvidas.

\begin{table}[H]
\centering\scriptsize\caption{Escores históricos da MIM da C1, sujeitos a reavaliação sob M1}\label{tab:mim-pontos}
\begingroup\setlength{\tabcolsep}{2pt}\renewcommand{\arraystretch}{1.25}
\begin{tabularx}{\textwidth}{L{3.8cm}*{13}{>{\centering\arraybackslash}X}}
\toprule\textbf{Direcionador}&\textbf{S01}&\textbf{S02}&\textbf{S03}&\textbf{S04}&\textbf{S05}&\textbf{S06}&\textbf{S07}&\textbf{S08}&\textbf{S09}&\textbf{S10}&\textbf{S11}&\textbf{S12}&\textbf{S13*}\\\midrule
Carry-over & 1 & 3 & 3 & 1 & 3 & 3 & 1 & 3 & 3 & 3 & 3 & 3 & 0 \\
Evolução tecnológica & 1 & 0 & 1 & 3 & 0 & 0 & 0 & 0 & 0 & 0 & 0 & 0 & 5 \\
Alterações planejadas de projeto & 0 & 0 & 0 & 0 & 0 & 0 & 0 & 0 & 0 & 0 & 0 & 0 & 0 \\
Especificação técnica & 5 & 3 & 5 & 3 & 1 & 1 & 1 & 1 & 1 & 0 & 1 & 3 & 3 \\
Estilo & 0 & 0 & 0 & 0 & 0 & 0 & 0 & 0 & 0 & 0 & 0 & 0 & 0 \\
Unidade comum & 3 & 5 & 3 & 1 & 5 & 5 & 3 & 5 & 5 & 5 & 5 & 3 & 1 \\
Processo e organização & 5 & 1 & 3 & 3 & 1 & 1 & 3 & 3 & 1 & 1 & 1 & 3 & 1 \\
Testes em separado & 5 & 3 & 3 & 5 & 1 & 3 & 5 & 3 & 3 & 3 & 3 & 3 & 5 \\
Compra de produto pronto & 3 & 3 & 1 & 1 & 0 & 0 & 1 & 0 & 3 & 0 & 0 & 1 & 5 \\
Manutenção e mantenabilidade & 5 & 5 & 3 & 3 & 1 & 1 & 5 & 1 & 3 & 1 & 1 & 3 & 3 \\
Atualização & 3 & 1 & 3 & 3 & 0 & 1 & 1 & 1 & 1 & 0 & 1 & 1 & 5 \\
Reciclagem & 5 & 3 & 3 & 3 & 1 & 1 & 5 & 3 & 3 & 0 & 1 & 3 & 3 \\
\midrule\textbf{Soma bruta}&36 & 27 & 28 & 26 & 13 & 16 & 25 & 20 & 23 & 13 & 16 & 23 & 31\\
\textbf{Posição na base}&1 & 3 & 2 & 4 & 11 & 9 & 5 & 8 & 6 & 11 & 9 & 6 & --\\
\bottomrule\end{tabularx}\endgroup
\fonte{Julgamentos anteriores à mudança para M1, preservados para rastreio. S13* não participa da base. As notas ainda não foram revalidadas para proteção direta.}
\end{table}

Na avaliação histórica, amortecimento (S01, 36) reuniu motivos fortes em fabricação, teste, manutenção e separação material. Ajuste (S03, 28) e imobilização (S02, 27) concentraram variação dimensional e reposição de desgaste. Inspeção (S04, 26) e retenção (S07, 25) exigiram combinar acesso e fechamento. Esses resultados orientam a reavaliação dos perfis sob M1; não são uma seleção das quatro maiores colunas nem importância funcional do cliente.

A planilha calcula somas e posições com empates, inclui OP-01 apenas pelo controle explícito e impede uma soma válida com nota ausente ou fora de 0/1/3/5. O registro conserva a nota histórica e sinaliza a reavaliação da fronteira M1; se a nota mudar, sinaliza também a revisão do julgamento. As somas continuam calculáveis como histórico, sem autorizar a transferência para a nova arquitetura. Não há agrupamento automático por pontuação.

Uma análise de sensibilidade substituiu 5 por 9 apenas para comparar a convenção da base com a escala consultada. As posições mudam: por exemplo, retenção passa de 25 para 37, e imobilização de 27 para 35. Os motivos para separar amortecimento, ajuste/imobilização e barreira/inspeção permanecem identificáveis; o corpo comum continua sustentado por interfaces compartilhadas. Esse teste não prova robustez a qualquer julgamento. A ordenação não deve ser usada como corte arbitrário do número de módulos.

\subsection{Arquitetura do Produto}
\label{sec:arquitetura-mim}

A análise anterior de perfis e vínculos físicos indicou quatro núcleos para C1: MD1 estrutura; MD2 ajuste e imobilização; MD3 amortecimento; MD4 contenção e inspeção. Eles permanecem hipóteses a reavaliar sob M1. MD6 é opção de evento, fora da referência. Uso e identidade do antigo MD5 foram integrados a MD1, pois cava e gravação pertencem ao corpo. A numeração de MD6 foi preservada para rastrear a opção anterior.

\begin{table}[H]
\centering\scriptsize\caption{Agrupamentos herdados a reavaliar para a proteção direta}\label{tab:mim-modulos}
\begin{tabularx}{\textwidth}{L{0.8cm}L{3cm}L{3.1cm}Y}
\toprule\textbf{ID}&\textbf{Núcleo}&\textbf{Portadores}&\textbf{Motivo e limite}\\\midrule
MD1 & Estrutura, pega e circulação & S05, S06, S08, S09, S10, S11, S12 & Predominam unidade comum e carry-over. Pega, carga, apoio, recolhimento e identidade compartilham o corpo; separar cada função aumentaria interfaces estruturais. Fechos podem ser peças de reposição internas, sem se tornarem automaticamente módulos. \\
MD2 & Ajuste e imobilização & S02, S03 & Cintas e guias atuam sobre a peça descoberta; coordenar curso, tensão, área de apoio e pressão na cerâmica. Reavaliar os motivos de agrupamento herdados. \\
MD3 & Amortecimento & S01 & Processo, teste em separado, manutenção e reciclagem são fortes. O calço tem material e vida próprios, sem exigir descarte do corpo. \\
MD4 & Contenção e inspeção & S04, S07 & Teste, manutenção, processo e fim de vida indicam separação da estrutura. O agrupamento entre visor e barreira é condicional ao fechamento contínuo e ao acesso real à cerâmica; seus perfis não são idênticos. \\
MD6* & Registro de evento opcional & S13 & Tecnologia, fornecedor, teste separado e atualização favorecem separação se OP-01 entrar no escopo. Alto escore não autoriza sua inclusão. \\
\bottomrule\end{tabularx}
\fonte{Análise própria dos perfis da MIM e dos portadores da C1; MD6* não incluído.}
\end{table}

As funções de pega, compressão, apoio, recolhimento e identidade permanecem ligadas ao corpo MD1. Fechos são partes locais de interface, sem serem contados automaticamente como módulos. S12 tem motivo de manutenção, mas separar o recolhimento como núcleo autônomo pode criar juntas no caminho de carga; a decisão permanece condicionada à estabilidade e à compressão. MD4 reúne perfis diferentes por exigir continuidade da barreira ao redor da região de inspeção. Essa integração só permanece se ambas as metas forem verificadas no mesmo conjunto.

Não se somaram escores por módulo para comparar grupos de tamanhos diferentes. A MIM não prova quantidade de módulos substituíveis. Para RP20, a hipótese é demonstrar a troca independente de MD2, MD3 e MD4. MD1 preserva a referência estrutural e a identidade da unidade; não entra nessa contagem. MD6 opcional também não foi usado para atingir três módulos.

\begin{table}[H]
\centering\scriptsize\caption{Interfaces e verificação da troca independente}\label{tab:mim-interfaces}
\begin{tabularx}{\textwidth}{L{1.1cm}L{3.7cm}YY}
\toprule\textbf{ID}&\textbf{Interface a definir}&\textbf{Troca a executar}&\textbf{Verificação após reposição}\\\midrule
I12&Base--ajuste: fixação, datum, curso e trava&Retirar/repor MD2 sem desmontar MD3/MD4&Amplitude sem troca no ajuste e deslocamento residual\\
I13&Base/ajuste--calço: assento, contato e retenção&Retirar/repor MD3 mantendo os outros núcleos&Resposta do calço e aceleração, dispersão e vida do conjunto\\
I14&Base/ajuste--barreira: fixação, vedação e acesso óptico&Retirar/repor MD4 preservando base, guias e calços&Área/defeitos verificáveis, retenção e retirada segura\\
I16*&Base--OP-01: montagem, posição, energia/leitura&Somente após inclusão no escopo&Calibração, benefício, vida e custo; sem diagnóstico de trinca\\
\bottomrule\end{tabularx}
\fonte{Plano de verificação proposto; resultados ainda não disponíveis. I16* está fora da C1 base.}
\end{table}

Executar trocas em conjunto descarregado, registrando acesso, ações, tempo e partes retiradas. Ajustar a faixa de tamanho permanece na mesma configuração: não se confunde com remover módulo por manutenção. Após cada troca, repetir as verificações afetadas e preservar a identidade. C2 e C3 exigem reavaliação da MIM para seus próprios portadores, particularmente rolamento, sensores e identificação; esta arquitetura não foi transportada automaticamente para essas concepções.


\subsubsection{Concepção 1}

A C1 busca simplicidade e reposição de desgaste. Propõe suporte polimérico, EPP, cintas e guias reguláveis na mesma configuração, com interface para carrinho. A peça entra descoberta após preparo no distribuidor; cintas e apoios precisam limitar a pressão de contato. O visor integra a barreira do próprio dispositivo, removível para manutenção. A versão base não inclui indicador. Área crítica e contenção precisam ser verificadas no mesmo conjunto.

\subsubsection{Concepção 2}

A C2 explora estrutura telescópica, interface resiliente em contato com a peça, fechos indexados, painel transparente na barreira e rolamento removível. Pressão nos apoios, limpeza, massa, inspeção e retorno são riscos a verificar. Registrador eletrônico fica fora da base e só pode ser estudado separadamente, com benefício, energia, memória, leitura e calibração definidos.

\subsubsection{Concepção 3}

A C3 explora rodízios com freio, apoios de descanso, painéis reguláveis com regiões transparentes e barreira fechada em torno da peça descoberta. Contato com a cerâmica, área crítica, escadas, desníveis e estabilidade com rodas precisam de avaliação. Indicador comercial fica fora da base e pode ser estudado como complemento calibrado.

\subsubsection{Concepção Final}

Não há concepção final selecionada. A primeira seleção priorizou a fronteira M1; a segunda Pugh compara C1/C2/C3 dentro dessa fronteira. C1 é datum, sem aprovação automática, e C2/C3 permanecem AV. Cada nota precisa de comparação na mesma carga, percurso e condições, com resultado ou estimativa identificada, justificativa da intensidade e limite ainda a confirmar. Dimensionar e analisar as propostas; executar ensaios quando a afirmação depender de desempenho físico. Todas as restrições continuam aplicáveis.

\subsection{Análise Técnica}

A análise técnica está planejada, sem resultados de ensaio. Deve avaliar transferência de choque/vibração, deslocamento, compressão, umidade, inspeção e retenção. No canteiro, separar empurrar/puxar, levantar/carregar em dupla e transpor desníveis. A referência de 150\,N do RP3 é de força de partida ao empurrar/puxar, não limite de massa ou de levantamento \citep{hsepushpull}.

A configuração carregada desfavorável deve ser avaliada para tombamento. A força é medida no piso de uso. O envelope inclui rodas, pegas e apoios na configuração de transporte. Sob M1, definir os assentos diretamente na cerâmica, limitar pressão e tensão e verificar que o caminho de compressão não concentra carga na louça. O preparo no distribuidor inclui retirar a embalagem, conferir a condição inicial, transferir a peça com segurança, identificar a unidade e encaminhar o material retirado. A embalagem de origem deixa de integrar o trecho carregado após esse preparo; o envelope do novo dispositivo continua sujeito a RP8 e RP14.

\subsection{Análise de Valor}

A análise de valor depende de funções e componentes definidos. Faltam custos de fabricação, aquisição, manutenção e retorno. As variantes devem demonstrar benefício da proteção e da inspeção frente aos custos e à complexidade. O registro opcional precisa de valor incremental; não é aprovado apenas por acrescentar informação.

\subsection{FMEA}

Falhas a investigar: deslocamento dentro da caixa; aceleração excessiva sem ativação do indicador; indicador ativado sem dano; inspeção obstruída; escape de fragmentos na abertura; fecho solto; tombamento/rolamento; perda de ajuste; e perda de identidade após troca. Esses itens orientam uma futura FMEA. Severidade, ocorrência, detecção e prioridade não foram pontuadas sem base.

\subsection{Análise de Viabilidade}

A unidade de comparação é uma entrega no fluxo estudado. O custo do dispositivo por ciclo é uma parcela. Na direção M1, acrescentam-se preparo no distribuidor e destino da embalagem retirada, além de retorno, manutenção, consumíveis e avarias residuais. O custo atribuível da embalagem de origem deve usar a mesma fronteira contábil da situação atual: retirar a caixa não demonstra deixar de pagá-la. Evitar dupla contagem entre aquisição, manuseio e destinação. O teto provisório de RP23 não autoriza omitir parcelas.

\begin{table}[H]
\caption{Cenário econômico histórico com caixa mantida, sem cotação das concepções}\centering\footnotesize
\begin{tabularx}{\textwidth}{L{6cm}L{3.2cm}Y}
\toprule\textbf{Parcela/comparação}&\textbf{Valor}&\textbf{Natureza}\\\midrule
Embalagem atual&R\$\,15/entrega&Hipótese de compra industrial\\
Avaria e reentrega&$0{,}02\times(500+120)$&Hipóteses a confirmar\\
Cenário atual&R\$\,27,40/entrega&Cálculo sobre as hipóteses\\
Dispositivo de R\$\,300 em 20 ciclos&R\$\,15/ciclo&Não é cotação da M3\\
Caixa mantida mais dispositivo&R\$\,30/entrega&Limite inferior sem retorno, manutenção, sinal ou avaria residual\\
Margem adicional com caixa mantida&R\$\,12,40/entrega&Limite otimista com avaria residual zero\\
\bottomrule\end{tabularx}
\fonte{Premissas da planilha. Taxa de avaria, preços e vida não foram medidos pela equipe.}
\end{table}

Esse cenário continua útil para M2/M3 e não foi reapresentado como custo de M1. Para a proteção direta, a planilha acrescenta um quadro próprio, com preparo, destinação e custo atribuível da embalagem. O cálculo completo depende dessas entradas e dos custos de retorno, manutenção, consumíveis e avaria residual. Aquisição de R\$\,300 e vida de vinte ciclos permanecem hipóteses de cenário; não são cotação nem vida medida. O resultado fica indisponível enquanto faltarem entradas. A escolha funcional não demonstra vantagem econômica.

\subsection{Modelo de Negócio}

ATO e produção própria permanecem premissas de planejamento. O modelo comercial precisa definir quem detém o dispositivo, como ocorre retorno, quem paga manutenção e como são tratadas perdas. Isso afeta custo e rastreio. Não se declarou demanda, preço de venda ou retorno financeiro validado.

\noindent\textbf{Encaminhamento.} Desenvolver a proteção direta da M1 e o preparo no distribuidor. Confirmar mix, contatos, inspeção, contenção e movimentação; medir a situação atual e obter custos comparáveis. Avançar os princípios e interfaces de C1/C2/C3, reavaliando os escores herdados da MIM. Com análises e ensaios necessários, justificar as notas da segunda Pugh e selecionar uma concepção que respeite as restrições. Reabrir a fronteira funcional se M1 se mostrar inviável. O detalhamento e os ensaios de confirmação seguem essa seleção. Planilha e relatório devem ser atualizados na mesma rodada.
